%=================================================
% CHAPITRE 3 : SPRINT 1
%=================================================

\chapter{Sprint 1 : Socle technique, Multi-Tenancy et isolation temporelle}

\section{Introduction}

Quand j'ai commencé à poser les premières briques du projet, ma priorité n'était pas de construire de jolis écrans, mais de m'assurer que les données soient parfaitement isolées. Dans un contexte universitaire, on manipule des informations sensibles (diplômes, contrats, salaires) et il est impensable que les données d'une faculté puissent être visibles par une autre, ou que les contrats de l'année passée se mélangent avec ceux de l'année en cours.

Ce premier sprint a donc servi à bâtir les fondations invisibles mais indispensables de SIRH-Doc : l'architecture multi-tenant et la gestion du contexte temporel par année universitaire.

\section{Objectifs et Backlog du sprint}

L'idée était d'arriver, à la fin de cette itération, à une application capable de reconnaître qui se connecte, pour quelle organisation il travaille, et sur quelle période académique il souhaite agir.

\begin{itemize}
    \item \textbf{Mise en place du Multi-Tenancy} : Capacité du backend à aiguiller les requêtes vers la bonne base de données.
    \item \textbf{Sécurité JWT} : Authentification robuste avec des jetons qui portent l'identité et le périmètre de l'utilisateur.
    \item \textbf{Contexte d'Année Universitaire} : Création d'un filtre global pour cloisonner les données RH par session.
\end{itemize}

\section{Choix d'architecture et défis rencontrés}

\subsection{La résolution dynamique du Tenant}

Au départ, j'avais envisagé une gestion classique via session utilisateur côté serveur, mais cette approche devenait vite compliquée avec les changements d'organisation. J'ai finalement opté pour un middleware de \textit{Tenant Resolution}. À chaque appel HTTP, le système regarde le jeton de l'utilisateur et configure instantanément la connexion SQL via une \texttt{ConnectionFactory}. C'est transparent pour le reste du code, ce qui m'a évité de passer l'identifiant de l'organisation dans chaque fonction.

\subsection{L'année universitaire comme contexte global}

C'est sans doute le point qui m'a demandé le plus de réflexion. On ne peut pas simplement mettre une date dans une table. Un vacataire peut avoir un contrat en 2024-2025 et un autre en 2025-2026. J'ai donc choisi de traiter l'année universitaire comme un \og état global \fg{} de l'application. Une fois sélectionnée côté frontend, elle est injectée dans toutes les requêtes SQL via Dapper pour filtrer les parcours et les positions administratives.

\section{Analyse et Conception}

\subsection{Diagramme de cas d'utilisation}

Le diagramme suivant montre comment le point d'entrée du système est verrouillé. Avant de pouvoir faire quoi que ce soit, l'utilisateur doit franchir l'étape de l'authentification et du choix de son contexte de travail (organisation et année).

\diagramplaceholder{sprint1_cas_utilisation}{Cas d'utilisation du Sprint 1}{Initialisation du contexte et sécurité}

\subsection{Diagramme de classes (Entités persistantes)}

Conformément à mon choix d'isoler les responsabilités, les classes persistantes de ce sprint se concentrent sur la gestion des accès et des structures de base. On y retrouve l'utilisateur, son organisation de rattachement et le référentiel des années académiques.

\diagramplaceholder{sprint1_classes}{Diagramme de classes - Sprint 1}{Entités de base pour la sécurité et le contexte}

\subsection{Diagrammes de séquence : Flux techniques du socle}

Les flux suivants illustrent les deux mécanismes fondamentaux de ce sprint : la sécurisation des accès et l'isolation automatique des données.

\subsubsection{Authentification et génération de jeton (JWT)}

Le flux de connexion illustre bien la complexité cachée derrière un simple bouton \og Valider \fg{}. Le backend ne se contente pas de vérifier un mot de passe ; il construit un jeton JWT contenant le périmètre de l'utilisateur (rôle et organisation) pour sécuriser les appels futurs.

\diagramplaceholder{sprint1_sequence_auth}{Séquence d'authentification}{Processus de login et génération du périmètre utilisateur}

\subsubsection{Résolution dynamique du Tenant (Middleware)}

Ce diagramme détaille le fonctionnement du middleware .NET. À chaque requête, le système intercepte le jeton JWT et les en-têtes HTTP pour identifier l'établissement et l'année universitaire. Cette information est injectée dans un contexte \textit{Scoped} pour filtrer automatiquement toutes les requêtes SQL.

\diagramplaceholder{sprint1_sequence_middleware}{Séquence de résolution de Tenant}{Interception et injection du contexte organisationnel}

\section{Réalisation technique}

\subsection{Backend : Sécurisation et Dapper}

Pour le stockage des mots de passe, je n'ai pris aucun risque : j'ai utilisé \textbf{BCrypt}. Côté accès aux données, \textbf{Dapper} s'est avéré être un excellent choix pour sa rapidité, surtout quand il s'agit d'injecter dynamiquement des clauses \texttt{WHERE} liées au TenantId ou à l'année universitaire.

\subsection{Frontend : Le Guard de sélection}

Côté Angular, j'ai mis en place un \textit{Guard} spécifique. Si un administrateur tente d'accéder à la gestion du personnel sans avoir choisi une année universitaire active, il est automatiquement renvoyé vers un écran de sélection. Cela garantit qu'aucune donnée n'est créée \og dans le vide \fg{}.

\section{Conclusion du sprint}

À la fin de ce sprint, j'avais un socle solide. L'application savait qui était connecté et dans quel cadre temporel il travaillait. C'était l'étape nécessaire avant de pouvoir attaquer la gestion réelle des dossiers RH dans le sprint suivant.
